% from ega2-1-fr.tex
\oldpage[II]{78}

homomorphismes de $\mathcal{O}_Y$-Modules quasi-cohérents aux homomorphismes
surjectifs. Il faut en effet montrer que si
$\mathcal{E}\to\mathcal{E}'$ est un $\mathcal{O}_Y$-homomorphisme surjectif,
le diagramme
\[
  \xymatrix{
    \mathbf{P}(\mathcal{E}')\times_Y\mathbf{P}(\mathcal{F})
      \ar[r]^{j\times1} \ar[d]_{\varsigma}
    & \mathbf{P}(\mathcal{E})\times_Y\mathbf{P}(\mathcal{F})
      \ar[d]^{\varsigma}
    \\\mathbf{P}(\mathcal{E}'\otimes\mathcal{F}) \ar[r]
    & \mathbf{P}(\mathcal{E}\otimes\mathcal{F})
  }
\]
est commutatif, $j$ désignant l'immersion fermée canonique
$\mathbf{P}(\mathcal{E}')\to\mathbf{P}(\mathcal{E})$. Posons
$P_1'=\mathbf{P}(\mathcal{E}')$ et gardons par ailleurs les notations de
(\hyperref[II.4.3.1-fr]{4.3.1})~; $j\times1$ est une immersion fermée
(\hyperref[I.4.3.1-fr]{I, 4.3.1}) et on a, à des isomorphismes près,
\[
  (j\times1)^*(\mathcal{O}_{P_1}(1)\otimes\mathcal{O}_{P_2}(1))
  =j^*(\mathcal{O}_{P_1}(1))\otimes\mathcal{O}_{P_2}(1)
  =\mathcal{O}_{P_1'}(1)\otimes\mathcal{O}_{P_2}(1)
\]
en vertu de (\hyperref[II.4.1.2.1-fr]{4.1.2.1}) et de
(\hyperref[I.9.1.5-fr]{I, 9.1.5})~; notre assertion résulte donc de
(\hyperref[II.4.2.8-fr]{4.2.8}) et
(\hyperref[II.4.2.9-fr]{4.2.9}).
\end{env}

\begin{env}[4.3.5]
\label{II.4.3.5-fr}
Avec les notations de (\hyperref[II.4.3.1-fr]{4.3.1}), soit
$\psi:Y'\to Y$ un morphisme, et posons
$\mathcal{E}'=\psi^*(\mathcal{E})$, $\mathcal{F}'=\psi^*(\mathcal{F})$~;
alors le morphisme de Segre
$\mathbf{P}(\mathcal{E}')\times\mathbf{P}(\mathcal{F}')
\to\mathbf{P}(\mathcal{E}'\otimes\mathcal{F}')$ s'identifie à
$\varsigma_{(Y')}$. En effet, gardant les notations de
(\hyperref[II.4.3.1-fr]{4.3.1}), posons par ailleurs
$P_1'=\mathbf{P}(\mathcal{E}')$, $P_2'=\mathbf{P}(\mathcal{F}')$~; on sait
(\hyperref[II.4.1.3.1-fr]{4.1.3.1}) que $P_i'$ s'identifie à $(P_i)_{(Y')}$
($i=1,2$), donc le morphisme structural $P_1'\times P_2'\to Y'$ s'identifie
à $r_{(Y')}$. D'autre part,
$\mathcal{E}'\otimes\mathcal{F}'$ s'identifie à
$\psi^*(\mathcal{E}\otimes\mathcal{F})$, donc
$\mathbf{P}(\mathcal{E}'\otimes\mathcal{F}')$ s'identifie à
$(\mathbf{P}(\mathcal{E}\otimes\mathcal{F}))_{(Y')}$
(\hyperref[II.4.1.3.1-fr]{4.1.3.1}). Enfin,
$\mathcal{O}_{P_1'}(1)\otimes_{Y'}\mathcal{O}_{P_2'}(1)=\mathcal{L}'$
s'identifie à $\mathcal{L}\otimes_Y\mathcal{O}_{Y'}$, en vertu de
(\hyperref[II.4.1.3.2-fr]{4.1.3.2}) et de
(\hyperref[I.9.1.11-fr]{I, 9.1.11}). L'homomorphisme canonique
$(r_{(Y')})^*(\mathcal{E}'\otimes\mathcal{F}')\to\mathcal{L}'$ s'identifie
alors à $s_{(Y')}$, et notre assertion résulte de
(\hyperref[II.4.2.10-fr]{4.2.10}).
\end{env}

\begin{remark}[4.3.6]
\label{II.4.3.6-fr}
Le préschéma somme de $\mathbf{P}(\mathcal{E})$ et
$\mathbf{P}(\mathcal{F})$ est de même canoniquement isomorphe à un
sous-préschéma fermé de $\mathbf{P}(\mathcal{E}\oplus\mathcal{F})$. En
effet, les homomorphismes surjectifs
$\mathcal{E}\oplus\mathcal{F}\to\mathcal{E}$,
$\mathcal{E}\oplus\mathcal{F}\to\mathcal{F}$ correspondent à des immersions
fermées
$\mathbf{P}(\mathcal{E})\to\mathbf{P}(\mathcal{E}\oplus\mathcal{F})$,
$\mathbf{P}(\mathcal{F})\to\mathbf{P}(\mathcal{E}\oplus\mathcal{F})$~;
tout revient à voir que les espaces sous-jacents des sous-préschémas fermés
de $\mathbf{P}(\mathcal{E}\oplus\mathcal{F})$ ainsi obtenus sont sans point
commun. La question étant locale sur $Y$, on peut adopter les notations de
(\hyperref[II.4.3.1-fr]{4.3.1})~; or $\mathbf{S}_n(E)$ et
$\mathbf{S}_n(F)$ s'identifient à des sous-modules de
$\mathbf{S}_n(E\oplus F)$ d'intersection réduite à $0$~; et si
$\mathfrak{p}$ est un idéal premier gradué de $\mathbf{S}(E)$ tel que
$\mathfrak{p}\cap\mathbf{S}_n(E)\neq\mathbf{S}_n(E)$ pour tout $n\geq0$,
il lui correspond dans $\mathbf{S}(E\oplus F)$ un idéal premier gradué dont
la trace sur $\mathbf{S}_n(E)$ est
$\mathfrak{p}\cap\mathbf{S}_n(E)$, mais qui contient $\mathbf{S}_+(F)$,
comme on le voit aussitôt~; deux points de
$\operatorname{Proj}(\mathbf{S}(E))$ et
$\operatorname{Proj}(\mathbf{S}(F))$ respectivement ne peuvent donc avoir
même image dans $\operatorname{Proj}(\mathbf{S}(E\oplus F))$.
\end{remark}

\subsection{Immersions dans les fibrés projectifs. Faisceaux très amples}
\label{subsection:II.4.4-fr}

\begin{proposition}[4.4.1]
\label{II.4.4.1-fr}
Soient $Y$ un schéma quasi-compact, ou un préschéma dont l'espace sous-jacent
est noethérien, $q:X\to Y$ un morphisme de type fini, $\mathcal{L}$ un
$\mathcal{O}_X$-Module inversible.
\begin{enumerate}
  \item[\rm{(i)}] Soient $\mathcal{S}$ une $\mathcal{O}_Y$-Algèbre graduée
  quasi-cohérente à degrés positifs,
  $\psi:q^*(\mathcal{S})\to\bigoplus_{n\geq0}\mathcal{L}^{\otimes n}$ un
  homomorphisme d'Algèbres graduées. Pour que $r_{\mathcal{L},\psi}$ soit
  partout défini et soit une immersion, il faut et

